\begin{exsol}{16}{Tadelis Exercise 4.7}{77--78}{Exam}{B3}
\asked Two candidates choose a positive ($P$), balanced ($B$) or negative ($N$) campaign; payoffs
are win probabilities given by the book's rules. (a) Model it as a normal-form game; (b) write
the matrix; (c) describe each stage of IESDS and say whether it gives a clear prediction.

\method Normal form (slides~13--14); IESDS (slide~29; Tadelis \S4.2).

\solution
\textbf{(a)} $I = \{1, 2\}$, $S_1 = S_2 = \{P, B, N\}$. Let $\pi(s_i, s_j)$ be the probability that
the candidate playing $s_i$ wins against $s_j$. The book's rules give
\[
  \pi(s,s) = 0.5,\quad \pi(P,B) = 0,\quad \pi(P,N) = 0.3,\quad \pi(N,B) = 0.6 .
\]
The other three entries follow because exactly one candidate wins:
$\pi(B,P) = 1$, $\pi(N,P) = 0.7$, $\pi(B,N) = 0.4$. Payoffs are $u_1(s_1,s_2) = \pi(s_1,s_2)$ and
$u_2(s_1,s_2) = 1 - u_1(s_1,s_2) = \pi(s_2,s_1)$. Player~2's payoff function is the same rule
applied with the roles swapped. It is the complement of player~1's because the probabilities of
the two candidates winning add up to one: the game is constant-sum and symmetric.

\textbf{(b)}
\begin{center}
\begin{tabular}{r|c|c|c|}
\multicolumn{1}{r}{1 $\backslash$ 2} & \multicolumn{1}{c}{$P$} & \multicolumn{1}{c}{$B$} & \multicolumn{1}{c}{$N$}\\ \cline{2-4}
$P$ & $0.5,\ 0.5$ & $0,\ 1$ & $0.3,\ 0.7$ \\ \cline{2-4}
$B$ & $1,\ 0$ & $0.5,\ 0.5$ & $0.4,\ 0.6$ \\ \cline{2-4}
$N$ & $0.7,\ 0.3$ & $0.6,\ 0.4$ & $0.5,\ 0.5$ \\ \cline{2-4}
\end{tabular}
\end{center}

\textbf{(c) IESDS.}
\begin{itemize}
  \item Round 1. Player~1's rows over columns $(P, B, N)$: $P = (0.5, 0, 0.3)$,
  $B = (1, 0.5, 0.4)$, $N = (0.7, 0.6, 0.5)$. $P$ is strictly dominated by $B$
  ($0.5 < 1$, $0 < 0.5$, $0.3 < 0.4$) and also by $N$. $B$ and $N$ are not ranked ($B$ is better
  against $P$, $N$ against $B$ and $N$). By symmetry the same holds for player~2.
  $S_i^1 = \{B, N\}$.
  \item Round 2. On $\{B, N\}$: $B = (0.5, 0.4)$, $N = (0.6, 0.5)$. $B$ is strictly dominated
  by $N$. The same holds for player~2. $S_i^2 = \{N\}$.
\end{itemize}
$S^\infty = \{(N, N)\}$: both run negative campaigns and each wins with probability $0.5$. The
prediction is clear (a unique profile) after two rounds.

\interp $B$ survives round~1 only because it crushes a positive campaign. Once each candidate
knows the other will not run a purely positive campaign, going fully negative wins. The
outcome has the same win probabilities as $(P,P)$, and in a constant-sum game every profile is
Pareto optimal \emph{for the candidates}. If voters dislike negative campaigning, the
equilibrium is a prisoner's-dilemma outcome for them: the candidates end up exactly where they
would have been with positive campaigns, but with a nastier campaign.

\verify \texttt{check\_A.py} builds the matrix from the four rules plus complementarity, checks
symmetry, that $P$ is dominated by both $B$ and $N$, and that IESDS removes $P$ (round~1) and
$B$ (round~2) for both players, leaving $(N, N)$.

\answer{(a) $S_i = \{P,B,N\}$, $u_1 = \pi(s_1,s_2)$ from the rules, $u_2 = 1 - u_1$ (the same
rule from candidate~2's side). (b) Matrix above. (c) Round~1 deletes $P$ (dominated by $B$ and
$N$); round~2 deletes $B$ (dominated by $N$); unique prediction $(N, N)$, each wins with
probability $0.5$.}
\end{exsol}
